%% \cvonly{\item ...} entries are left out of the resume's selected list.
\cvsection{\resumeonly{Selected Publications}\cvonly{Publications}}
141 citations (\link{https://scholar.google.com/citations?user=bgnVj8wAAAAJ}{Google Scholar}); research areas: \textbf{machine learning, representational drift, and whole-brain modeling}.
\par\addvspace{4pt}\textbf{Peer-reviewed journal articles}
\begin{itemize}
  \item \me{Wang, Y.}, Lu, Z., \& Lakshminarasimhan, K. (2026). Coordinated drift of image representations in human EEG. \status{under review}
  \item Lu, Z., \me{Wang, Y.}, \& Golomb, J. D. (2026). Achieving more human brain-like vision via human EEG representational alignment. \venue{Communications Biology}, 9, 463. \link{https://doi.org/10.1038/s42003-026-09685-w}{doi:10.1038/s42003-026-09685-w}
  \item Lu, Z., \& \me{Wang, Y.} (2025). Teaching CORnet human fMRI representations for enhanced model--brain alignment. \venue{Cognitive Neurodynamics}, 19(1), 61.
  \cvonly{%
  \item Li, Z., Huang, G., Li, Z., Li, S., \me{Wang, Y.}, Zhao, J., \ldots\ \& Zou, L. (2020). Chemosensory anhedonia in patients with schizophrenia and individuals with schizotypy: A questionnaire study. \venue{Frontiers in Psychiatry}, 11, 481.
  \item Zhao, J., Gao, Z.*, Li, Y.*, \me{Wang, Y.*}, Zhang, X., \& Zou, L. (2019). The food neophobia scale (FNS): Exploration and confirmation of factor structure in a healthy Chinese sample. \venue{Food Quality and Preference}, 79, 103791.
  \item Zhao, J.*, \me{Wang, Y.*}, Ma, Q., Zhao, J., Zhang, X., \& Zou, L. (2019). The Chemosensory Pleasure Scale: A new assessment for measuring hedonic smell and taste capacities. \venue{Chemical Senses}, 44(7), 457--464.
  \item \me{Wang, Y.}, \& Zou, L. (2018). The neural mechanisms of social pain. \venue{Progress in Biochemistry and Biophysics}, 45(7), 714--722. [in Chinese]}
\end{itemize}
\cvonly{\par\addvspace{2pt}*Equal contribution.}
\par\addvspace{4pt}\textbf{Conference papers and presentations}
\begin{itemize}
  \item \me{Wang, Y.}, \& Lakshminarasimhan, K. (2026). An efficient coding account of representational drift. \venue{International Conference on Learning Representations (ICLR 2027)}. \status{under review}
  \item Hosseinpourkhoshkbari, R., Huang, W., \me{Wang, Y.}, Helabad, A., Chen, Z., \& Golden, R. (2026). Human-likeness is multi-dimensional: Evaluating LLM-generated virtual patients against simulated patient dialogue. \venue{Conference on Language Modeling (COLM 2026)}. \status{accepted}
  \item \me{Wang, Y.}, \& Lakshminarasimhan, K. (2026). The geometry of drifting spatial representations. \venue{Cognitive Computational Neuroscience (CCN 2026)}. \status{accepted poster}
  \cvonly{%
  \item \me{Wang, Y.}, \& Lakshminarasimhan, K. (2026). When drift helps or hurts: Heterogenous effects of representational drift on decoding. \venue{Summit in BrainHealth, UT Dallas}.}
  \item Lu, Z., \me{Wang, Y.}, \& Golomb, J. D. (2024). ReAlnet: Achieving more human brain-like vision via human neural representational alignment. \venue{ICLR 2024 Workshop}. \status{accepted}
\end{itemize}
